%%%%%%%%%%  scan idx 218  |  folio 205  %%%%%%%%%%
qui est un
épimorphisme, comme on le constate aussitôt. Dans le cas où $Y=S\times\R$, le faisceau $\underline{\Homop}(X,Y)$ s'identifie canoniquement à $\cO_X$, et l'épimorphisme (2.1.1) à l'épimorphisme
\[
  i_s^{-1}(\cO_X)\longrightarrow i_s^*(\cO_X)=\cO_{X_s}
\]
considéré en (1.1).

\textbf{Proposition}\footnote{Communiquée au rédacteur par A. Douady.} \textbf{2.2.} Soient $X,Y$ des $S$-variétés mixtes, et $s\in S$. On suppose que la fibre $X_s$ est paracompacte et que $s$ admet un système fondamental de voisinages paracompacts. Soient $A$ un fermé de $X_s$, et $f\in H^0(A,\underline{\Homop}(X_s,Y_s))$.
Le sous-faisceau $F$ de $i_s^{-1}\underline{\Homop}(X,Y)$ formé des germes «~qui prolongent $f$~» (i.e. dont l'image par (2.1.1) coïncide avec $f$ au-dessus de $A$) est mou.

\emph{Preuve.} Soit $X'$ un voisinage ouvert de $x$ dans $X$, isomorphe (comme variété mixte) à $S'\times U$, où $S'$ est un voisinage paracompact de $s$ et $U$ un ouvert relativement compact de $\R^m$, et tel que l'image par $f$ de $X'_s\cap A$ soit contenue dans un ouvert de $Y$ isomorphe à $S'\times V$, ou $V$ est un ouvert convexe de $\R^n$ contenant l'origine. D'après (T.F. II 3.4.1), il suffit de montrer que toute section de $F$ au-dessus d'un fermé $A_0$ de $X_s$ contenu dans $X'_s$ se prolonge à $X'_s$. Or $X'_s$ ($\simeq\{s\}\times U$) est relativement compact, de sorte qu'un fermé de $X_s$ contenu dans $X'_s$ est compact. Identifions $X'$ (resp. $Y'$) à $S'\times U$ (resp. $S'\times V$), posons $A'=A\cap X'_s$, et soit $A_0$ un compact de $X'_s=\{s\}\times U$. Comme $A'$ est fermé dans $\{s\}\times U$, qui est paracompact, on a ([T.F.]~II 3.3.1)\footnote{Un fermé d'un espace paracompact admet un système fondamental de voisinages paracompacts.}
\[
  H^0\bigl(A',\underline{\Homop}(X_s,Y_s)\bigr)
  =\varinjlim_{U'\supset A'}\Homop(U',Y'_s),
\]
suivant les ouverts $U'$ de $X'_s$ contenant $A'$. De même, comme $A_0$ est fermé dans $S'\times \overline U$, qui est paracompact (produit d'un paracompact et d'un compact), on a
\[
  H^0\bigl(A_0,i_s^{-1}\underline{\Homop}(X,Y)\bigr)
  =\mathop{\underrightarrow{\mathrm{Lim}}}\limits_{U_0\supset A_0}\Homop(U_0,Y'),
\]
suivant les ouverts $U_0$ de $X'$ contenant $A_0$. La proposition sera donc
consé%
